\section{Pressure from weighted source layers}\label{sec:pressure}

We retain the sources, potentials, interaction measure, and cutoffs
defined above.  In particular, $n\ge3$, $A,B>-2$, $p,q>0$, and $pq>1$.
All constants in this section are uniform over admissible solution
pairs.  A subscript records permitted dependence on the auxiliary
parameters; constants are always independent of the direction $e$.

\begin{proposition}[Localized pressure]\label{prop:pressure}
Let $(u,v)$ be a positive solution pair as in Theorem~\ref{thm:main},
with $n\ge3$, $A,B>-2$, and $pq>1$.  For every $\eps>0$ there is a
constant $C_\eps$, depending only on the equation parameters and
$\eps$, such that
\begin{equation}\label{eq:pressure}
 aW_f+bW_g
 \le(n-2)\int_{\R^n\times\R^n}w(x)D(x,y)\,\dd\pi(x,y)
       +\eps E+C_\eps.
\end{equation}
\end{proposition}

The proof must accommodate the spatial factors at source-layer
endpoints.  We localize to short intervals, where those factors vary
by at most $C\delta d(x)$, and use caps of angular radius proportional
to $d(x)$.  The factor $d$ converts both errors into multiples of
$E$, by $dw\le C\eta$ from \eqref{eq:cutoff-bounds}.

\subsection{Angular localization}

Write $\mathbb E_e$ and $\mathbb P_e$ for expectation and probability
under normalized surface measure on $\Sph$.  Fix
$0<\kappa<1$ and $0<\delta<1/4$, and put
\[
 \ell(x)=\delta d(x)^m\qquad(d(x)>0).
\]
Choose a fixed nonnegative smooth function $\psi$ on $[0,\infty)$,
equal to one near zero, supported in $[0,1]$, and bounded by one.
For $0<l\le1$ and any $\theta\in\Sph$, set
\[
 Z(l)=\mathbb E_e\psi(|e-\theta|/l).
\]
Rotational invariance makes this independent of $\theta$.  Chordal
polar coordinates give
\[
 Z(l)=\frac{|\mathbb S^{n-2}|}{|\Sph|}
       l^{n-1}\int_0^1\psi(t)t^{n-2}
                  (1-l^2t^2/4)^{(n-3)/2}\,\dd t
\]
and therefore
$c l^{n-1}\le Z(l)\le C l^{n-1}$ and
$|Z'(l)|\le C l^{n-2}$.

Define the cap weight by
\begin{equation}\label{eq:cap-weight}
 h_e(x)=w(x)\frac{\psi(|e-\omega_x|/(\kappa d(x)))}{Z(\kappa d(x))}
 \quad(1<|x|<2),\qquad h_e(x)=0\quad\hbox{otherwise}.
\end{equation}
Its exact normalization is $\mathbb E_e h_e(x)=w(x)$.  Differentiation,
using \eqref{eq:cutoff-bounds}, gives
\begin{equation}\label{eq:cap-bounds}
 0\le h_e(x)\le C_\kappa d(x)^{N-n},\qquad
 |\nabla h_e(x)|\le C_\kappa d(x)^{N-n-1}
 \quad(1<|x|<2).
\end{equation}
To check the second estimate, derivatives of the cap profile cost at
most $C_\kappa/d$, and differentiation of $Z(\kappa d)^{-1}$ costs
the same relative factor.  The profile is constant near its center,
so the norm $|e-\omega_x|$ causes no differentiability problem there.
The weight extends continuously by zero across both boundary spheres:
$w$ vanishes quadratically at $|x|=1$, and $N-n>1$ at $|x|=2$.

As in Lemma~\ref{lem:localization}, call a pair close if
$|x-y|\le\ell(z)$ for at least one endpoint $z\in\{x,y\}$ with
$d(z)>0$.  On its segment, $d$ is comparable to $d(z)$ because $d$
is $1$-Lipschitz and $\ell(z)\le\delta d(z)$.  Integrating the
derivative bound across the pieces of this segment proves
\begin{equation}\label{eq:cap-oscillation}
 |h_e(x)-h_e(y)|
 \le C_\kappa\delta d(z)^{N-n-1+m}
 \le C_\kappa\delta d(z)^N,
\end{equation}
where $N-n-1+m=N+2$.  The interaction marginals bound the integral
over close pairs by $C_\kappa\delta E$.
On the complementary pairs, \eqref{eq:cap-bounds} and the tail bound
\eqref{eq:uniform} give, for example, the bound
\[
 C_\kappa\int_{B_2}f(x)d(x)^{N-n}\ell(x)^{2-n}\,\dd x
 \le C_{\kappa,\delta},
\]
since $N-n-m(n-2)=3n+8>0$; the other endpoint is identical with
$g$ in place of $f$.  We have proved
\begin{equation}\label{eq:cap-mismatch}
 \int|h_e(x)-h_e(y)|\,\dd\pi(x,y)
 \le C_\kappa\delta E+C_{\kappa,\delta}.
\end{equation}
Moreover, on the support of $h_e(x)$,
\[
 |D_e(x,y)-D(x,y)|\le2|e-\omega_x|\le2\kappa d(x).
\]
Averaging and using $dw\le C\eta$ gives
\begin{equation}\label{eq:cap-pressure}
 \mathbb E_e\int h_e(x)D_e(x,y)\,\dd\pi(x,y)
 \le\int w(x)D(x,y)\,\dd\pi(x,y)+C\kappa E.
\end{equation}
Here the constant $C$ is independent of $\kappa,\delta$.

\subsection{Measurable intervals and their loss}

For $0<t<f(x)$, let $I$ be the connected component containing $x$
of the source layer $\{f>t\}$ on the line $x+\R e$.  Assign it
the weight
\begin{equation}\label{eq:component-weight}
 \chi_I=
 \begin{cases}
  \displaystyle\inf_I h_e,
   &I\text{ bounded and }|I|\le2\delta\inf_I d,\\
  0,&\text{otherwise}.
 \end{cases}
\end{equation}
Set $\chi_f(e,x,t)=\chi_I$.  Define $\chi_g(e,y,t')=\chi_J$
in the same way from $\{g>t'\}$.  In particular,
$0\le\chi_f(e,x,t)\le h_e(x)$ and
$0\le\chi_g(e,y,t')\le h_e(y)$.

All these functions, as well as finite endpoints, are jointly
measurable.  Indeed, with the convention $f(0)=0$, the set
$\{(x,t):x\ne0,\ 0<t<f(x)\}$ is open.  The forward exit distance
of its line component exceeds $r\ge0$ exactly when the compact
segment $\{x+se:0\le s\le r\}$ is contained in the layer.  This
is an open condition in $(e,x,t)$; the backward exit distance is
treated in the same way.  Their sum, including the value $+\infty$,
is the measurable component length.  The infima in
\eqref{eq:component-weight} can be taken over rational offsets
between these two exits, since $d$ and $h_e$ are continuous.
This proves the assertion, including measurability in $e$, without
a choice of coordinates on the varying spaces $e^\perp$.

For each direction define the nonnegative losses
\begin{align*}
 L_f(e)&=\int_{\R^n}u(x)\int_0^{f(x)}
                 (h_e(x)-\chi_f(e,x,t))\,\dd t\,\dd x,\\
 L_g(e)&=\int_{\R^n}v(y)\int_0^{g(y)}
                 (h_e(y)-\chi_g(e,y,t'))\,\dd t'\,\dd y.
\end{align*}
They are finite for each fixed solution and direction, since $h_e$
is bounded and supported on the compact annulus $1\le|x|\le2$.
This finiteness does not assert a uniform bound on those marginals.

\begin{lemma}[Averaged layer loss]\label{lem:layer-loss}
For the cap weights \eqref{eq:cap-weight} and component weights
\eqref{eq:component-weight}, under the hypotheses of
Proposition~\ref{prop:pressure}, every $P\ge1$ satisfies
\begin{equation}\label{eq:layer-loss}
 \mathbb E_e\bigl(L_f(e)+L_g(e)\bigr)
 \le\bigl(C_\kappa\delta+C_{\kappa,\delta}P^{-1}\bigr)E+CP.
\end{equation}
The constants in the relative coefficient are independent of $P$,
and the last constant is independent of $\kappa,\delta,P$.
\end{lemma}

\begin{proof}
We prove the assertion for $f,u$; the same argument applies to $g,v$.
Only $1<|x|<2$ contributes.  Fix such an $x$ and write $d=d(x)$
and $\ell=\delta d^m$.
If the component $I$ at level $t$ has $|I|\le\ell$, then
\[
 \inf_I d\ge d-\ell\ge(1-\delta)d,
 \qquad |I|\le\delta d\le2\delta\inf_I d.
\]
It therefore passes the length condition in
\eqref{eq:component-weight}.  Estimate~\eqref{eq:cap-oscillation}
shows that its entire loss of weight is at most
$C_\kappa\delta d^N$.

If $|I|>\ell$, at least one of its two portions from $x$ contains
a segment of length $\ell/2$.  Let $F$ be the set of directions
whose forward segment of that length lies in $\{f>t\}$.  Polar
integration centered at $x$ gives
\[
 \int_{B_{\ell/2}(x)}f
 \ge \frac{|\Sph|}{n}\,t(\ell/2)^n\mathbb P_e(F).
\]
The long-interval event is contained in $F\cup(-F)$.  Antipodal
symmetry, followed by the source bound in \eqref{eq:uniform}, yields
\begin{equation}\label{eq:long-interval}
 \mathbb P_e\{|I|>\ell\}
 \le \frac{C}{t\ell^n}\int_{B_\ell(x)}f
 \le\frac{C}{t\ell^n}.
\end{equation}
Here $B_\ell(x)\subset B_4$; only its containment in a fixed
centered ball is used.

For $t\ge P d^{-k}$, the expected loss on long intervals is bounded
using the supremum of $h_e(x)$ times the probability in
\eqref{eq:long-interval}.  No independence is needed.  The result is
\[
 C_\kappa d^{N-n}\,C\delta^{-n}P^{-1}d^{k-mn}
   =C_{\kappa,\delta}P^{-1}d^N,
\]
because $k=n(m+1)$.  Integrating the short- and long-interval
estimates over these high levels, and then against $u(x)\,\dd x$,
gives at most
$\bigl(C_\kappa\delta+C_{\kappa,\delta}P^{-1}\bigr)E_f$.

For the remaining levels use the exact cap normalization:
\[
 \int_0^{\min\{f(x),Pd^{-k}\}}
        \mathbb E_e(h_e(x)-\chi_f(e,x,t))\,\dd t
 \le Pd^{-k}w(x)\le CPd^{N-1-k}\le CP.
\]
The last exponent equals $n+1$.  Integration against $u$ is bounded
by $CP$ by \eqref{eq:uniform}.  Repeating the proof for $g,v$
proves \eqref{eq:layer-loss}.
\end{proof}

\subsection{Weighted endpoints and nonnegative integration}

Suppose $\chi_I>0$ for an $f$-layer at level $t$.  Then $I$ is
bounded and $h_e\ge\chi_I>0$ on it.  Continuity gives the same
inequality at its two endpoints $x_-,x_+$, so they lie in the
annulus $1<|x|<2$.  Continuity of $f$ there and maximality of the
component imply
\[
 f(x_-)=f(x_+)=t,\qquad
 u(x_\pm)=t^{1/q}|x_\pm|^{-B/q}.
\]
For any $x\in I$, \eqref{eq:component-weight} gives
$|x_\pm-x|\le2\delta d(x)$.  The positive radial function
$|x|^{-B/q}$ has a bounded relative Lipschitz constant on
$1\le|x|\le2$.  Thus, with
$\overline u_I=(u(x_-)+u(x_+))/2$, we have
\begin{equation}\label{eq:weighted-endpoints}
 \left|\overline u_I-(t/|x|^B)^{1/q}\right|
 \le C\delta d(x)(t/|x|^B)^{1/q}
 \le C\delta d(x)u(x).
\end{equation}
The same statement holds for a positive-weight $g$-component $J$,
with endpoint average $\overline v_J$ and ideal value
$(t'/|y|^A)^{1/p}$.  When a component weight is zero, every product
involving its endpoint average is defined as zero without selecting
or evaluating endpoints.

Fix a direction $e$ and levels $t,t'>0$.  On two distinct parallel
lines $x=z+se$ and $y=z'+s'e$, apply
Lemma~\ref{lem:interval-pair} to each pair of layer components $I,J$
with $c_I=\chi_I$ and $c_J=\chi_J$.  Since
$\min(\chi_I,\chi_J)\le h_e(x)$, it gives
\begin{equation}\label{eq:layer-pair}
\begin{split}
 \int_{I\times J}(\chi_I+\chi_J)\mathcal K
 \le\int_{I\times J}\Bigl[
  &\chi_I\overline{\mathcal K}_I+\chi_J\overline{\mathcal K}_J\\
  &+\bigl((n-2)h_e(x)D_e(x,y)+|\chi_I-\chi_J|\bigr)
       \mathcal K\Bigr],
\end{split}
\end{equation}
where the measures are $\dd s\,\dd s'$.  At this stage all terms
are nonnegative extended integrals.

Each open line layer has at most countably many components.  Sum
\eqref{eq:layer-pair} over the two sets of components, and integrate
in $z,z'\in e^\perp$ and $t,t'>0$.  The summed functions are the
jointly measurable functions described above; hence Tonelli's
Theorem applies.  The coincident-line set $z=z'$ is null and may be
omitted.  By the potential identities \eqref{eq:potentials}, the two
terms on the left become
\begin{align*}
 U_f(e)&=\int\int_0^{f(x)}\chi_f(e,x,t)u(x)\,\dd t\,\dd x,\\
 U_g(e)&=\int\int_0^{g(y)}\chi_g(e,y,t')v(y)\,\dd t'\,\dd y.
\end{align*}
Similarly, the endpoint terms become
\begin{align*}
 V_f(e)&=\int\int_0^{f(x)}\chi_f(e,x,t)\overline u_I\,\dd t\,\dd x,\\
 V_g(e)&=\int\int_0^{g(y)}\chi_g(e,y,t')\overline v_J\,\dd t'\,\dd y.
\end{align*}
For example, integrating $\overline{\mathcal K}_I$ over all
$(y,t')$ layers recovers the average of $K*g$ at $x_-,x_+$.
For each selected endpoint the omitted $y$-points comprise one
affine line, which is Lebesgue null in $\R^n$, so the recovered
potential is exactly $\overline u_I$.

We now verify finiteness before any subtraction.  Set
\[
 H_f(e)=\int h_e(x)f(x)u(x)\,\dd x,
 \qquad H_g(e)=\int h_e(y)g(y)v(y)\,\dd y.
\]
Both are finite for every fixed solution and direction because the
cap weights are bounded and supported on a compact annulus.
We have $U_f\le H_f$, $U_g\le H_g$, while
\eqref{eq:weighted-endpoints} gives $V_f\le C H_f$ and
$V_g\le C H_g$.  The integrated pressure term is at most
$(n-2)H_f$.  Finally,
\[
 |\chi_f(e,x,t)-\chi_g(e,y,t')|\le h_e(x)+h_e(y),
\]
so its integrated contribution is at most $H_f+H_g$, again by
Tonelli and the two potential marginals.  All the integrated terms
are therefore finite, even when $n=3$.  In particular, we have
never subtracted the possibly infinite integrals over an unbounded
line component.

It is now legitimate to subtract $V_f+V_g$.  The sharper mismatch
bound
\begin{align*}
 |\chi_f(e,x,t)-\chi_g(e,y,t')|
 \le{}&|h_e(x)-h_e(y)|\\
      &+(h_e(x)-\chi_f(e,x,t))
        +(h_e(y)-\chi_g(e,y,t'))
\end{align*}
shows that the resulting finite quantity
\[
\begin{split}
 S_e={}&\int\int_0^{f(x)}\chi_f(e,x,t)(u(x)-\overline u_I)\,\dd t\,\dd x\\
   &+\int\int_0^{g(y)}\chi_g(e,y,t')(v(y)-\overline v_J)\,\dd t'\,\dd y
\end{split}
\]
satisfies
\begin{equation}\label{eq:integrated-pressure}
 S_e\le(n-2)\int h_e(x)D_e(x,y)\,\dd\pi(x,y)
       +\int|h_e(x)-h_e(y)|\,\dd\pi(x,y)+L_f(e)+L_g(e).
\end{equation}
The quantity $S_e$ need not be nonnegative; only the preceding
unsubtracted expressions were treated with Tonelli's Theorem.

\subsection{Recovery of the cutoff pressure}

The exact weighted layer identity is
\begin{equation}\label{eq:weighted-layer-identity}
 \int_0^{f(x)}\left[u(x)-(t/|x|^B)^{1/q}\right]\,\dd t
 =\frac{f(x)u(x)}{q+1}=a f(x)u(x)\qquad(x\ne0).
\end{equation}
Indeed, the integral of $(t/|x|^B)^{1/q}$ equals
$qf(x)u(x)/(q+1)$.  The analogous identity for $g$ has coefficient
$b=1/(p+1)$.
To compare \eqref{eq:weighted-layer-identity} with $S_e$, set
$\tau=(t/|x|^B)^{1/q}$.  Since $0\le\tau\le u(x)$,
\begin{align*}
 &h_e(x)(u(x)-\tau)-\chi_f(e,x,t)(u(x)-\overline u_I)\\
 &\quad=(h_e(x)-\chi_f(e,x,t))(u(x)-\tau)
                  +\chi_f(e,x,t)(\overline u_I-\tau)\\
 &\quad\le(h_e(x)-\chi_f(e,x,t))u(x)
                  +C\delta d(x)h_e(x)u(x),
\end{align*}
by \eqref{eq:weighted-endpoints}; the zero-weight case has the same
bound under the convention already stated.  Integrating and treating
the $g$-layer in the same way gives
\begin{equation}\label{eq:pressure-recovery}
\begin{split}
 aH_f(e)+bH_g(e)
 \le{}&S_e+L_f(e)+L_g(e)\\
       &+C\delta\int d(x)h_e(x)f(x)u(x)\,\dd x\\
       &+C\delta\int d(y)h_e(y)g(y)v(y)\,\dd y.
\end{split}
\end{equation}

\begin{proof}[Proof of Proposition~\ref{prop:pressure}]
Combine \eqref{eq:integrated-pressure} and
\eqref{eq:pressure-recovery}, and average over $e$.  This incurs
two copies of $L_f+L_g$.  The identities
$\mathbb E_e h_e=w$ and $dw\le C\eta$, together with
\eqref{eq:cap-mismatch}, \eqref{eq:cap-pressure}, and
Lemma~\ref{lem:layer-loss}, give
\[
\begin{split}
 aW_f+bW_g\le{}&(n-2)\int w(x)D(x,y)\,\dd\pi(x,y)\\
 &+\bigl(C\kappa+C_\kappa\delta
                    +C_{\kappa,\delta}P^{-1}\bigr)E
   +C_{\kappa,\delta,P}.
\end{split}
\]
The constant multiplying $\kappa$ is independent of all three
auxiliary parameters; the one multiplying $\delta$ is independent
of $\delta,P$; and the one multiplying $P^{-1}$ is independent of
$P$.  First choose $\kappa$ sufficiently small, next $\delta$
sufficiently small, and finally $P$ sufficiently large, so that
the relative coefficient is at most $\eps$.  The remaining
additive constant is finite and uniform over the solution pair.
This proves \eqref{eq:pressure}.
\end{proof}
